\DiaryEntry{Concurrent Events}{2016-08-24}{Stochastic}


Assume exactly two events occur in a timespan of \(T\). Each event lasts a time of \(d\). What is the probability that the two events overlap in time?

The Figure below shows the two events on the x and y axis,respectively. The shaded region is when the two events overlap.

\begin{figure}[H]
\centering
\includegraphics{images/concurrent_events.png}
\caption{Page1}
\end{figure}

In order to calculate the probability, we need to calculate areas:

The triangle below the shaded region has area \(\frac{1}{2}(T-d)^2\) and this is the same area as the triangle above. Therefore, the area of the shaded region is

\[
A = T^2 - (T-d)^2
\]

and the probability of an event overlap \(P\) is therefore

\[
P = \frac{T^2 - (T-d)^2}{T^2} = 1 - \left( \frac{T-d}{T}\right)^2
\]

In the extreme case of \(d=T\), we have \(P=1\) and for \(d=0\) we obtain \(P=0\). Both of these cases make intuitively sense.

In the case of \(d\) begin small, we can approximate \(P\) as follows

\[
P = 1 -  \frac{T^2 - 2dT + d^2}{T^2} = 1 - \left( 1 - 2\frac{d}{T} + \frac{d^2}{T^2} \right) \approx 2 \frac{d}{T}
\]
